Let @@M0@@ be a Chow presentation: @@M1@@ is projective over @@M2@@; and @@M3@@ is surjective. If we grant (a), (b), (c) and (iii) for @@M4@@, then they follow for @@M5@@ by 4.3 and 3.5. We may therefore suppose that @@M6@@ is projective over @@M7@@.

Let @@M8@@ be ample for @@M9@@. If, for an @@M10@@, @@M11@@ is ample, then @@M12@@ is ample for every generalization @@M13@@ of @@M14@@ (EGA III 4.7.1). By the equivalence of 4.6 (i) and (iv), we have therefore established (a).

For (b), (c) and 4.7 (iii), we may suppose that @@M15@@ is integral, with @@M16@@ the generic point in (b) and (c) (For (b) and (c), this is clear, and for (iii), it follows from 3.4 (i) applied to the morphism @@M17@@, where the @@M18@@ are the irreducible components of @@M19@@, equipped with the induced reduced structure).

There exist an integral scheme @@M20@@ and an f.p.p.f. morphism @@M21@@ such that every irreducible component @@M22@@ of @@M23@@, equipped with the induced reduced structure, is integral over @@M24@@ (XII 1.3). To prove (b) and (iii), it suffices to do so for @@M25@@, (in (b), by descent; in (iii), by noetherian induction and 3.4 (i)).

Finally, consider the canonical surjective morphism @@M26@@. Then (b) and (iii) for @@M27@@ follow from (c) and (iii) for the @@M28@@, by 4.3, 3.5 and 3.7.

We have therefore reduced to proving (c) and (iii), assuming @@M29@@ integral with generic point @@M30@@, and @@M31@@ projective and integral over @@M32@@. Let @@M33@@ be a very ample sheaf for @@M34@@. Then, for every @@M35@@, @@M36@@ and @@M37@@ do not depend on @@M38@@ (EGA III 7.9.4). Thus (c) follows from 4.6 (i) @@M39@@ (v), which also implies that the family of classes of invertible sheaves @@M40@@-equivalent to @@M41@@ on the fibers of @@M42@@ has just
